\begin{textsample}{POS Dim 1 – 2021\_c4 – Score -4.00 – 2021\_c4/enade\_2021\_ci\_bio\_d313503e1f.txt}  \label{ex:f1_pos_377}
CAPÍTULO 4 PERCEPÇÃO DA PROVA Nas análises feitas neste capítulo, constam as percepções dos concluintes da Área de Ciências Biológicas ( Licenciatura e Bacharelado ) sobre a prova aplicada no Enade/2021.
Essas percepções foram mensuradas por meio de nove questões que avaliaram desde o grau de dificuldade da prova até o tempo gasto para concluí -la.
As percepções sobre a prova foram relacionadas ao desempenho dos estudantes, à Grande Região de funcionamento do curso, à Categoria Administrativa e ao tipo de Organização Acadêmica da IES.
O questionário de percepção da prova encontra -se ao final dos Anexos VIIIA ( Licenciatura ) e VIIIB ( Bacharelado ), onde está a reprodução do exame.
O desempenho dos estudantes foi classificado em quatro quartos.
Para tanto, esse desempenho foi ordenado de forma ascendente.
O percentil 25, P25, também conhecido como primeiro quartil, é a nota de desempenho que deixa um quarto ( 25 \% ) dos valores observados abaixo e três quartos acima.
Na Figura 1, apresenta -se uma ilustração desse conceito.
O quarto inferior de desempenho é composto pelas notas abaixo do primeiro quartil.
Já o percentil 75, P75, também conhecido como terceiro quartil, é o valor para o qual há três quartos ( 75 \% ) dos dados abaixo e um quarto acima dele.
O quarto superior de desempenho é composto pelas notas iguais ou acima do terceiro quartil.
O percentil 50, P50, também conhecido como mediana, é o valor que divide as notas em dois conjuntos de igual tamanho.
O segundo quarto inclui os valores entre o primeiro quartil ( P25 ) e a mediana.
O terceiro quarto contém os valores entre a mediana ( P50 ) e o terceiro quartil ( P75 ).
Vale ressaltar que percentis, quartis e medianas são pontos que, não obrigatoriamente, pertencem ao conjunto original de dados, ao passo que os quartos são subconjuntos dos dados originais.
A seguir, serão apresentados gráficos com resultados selecionados, relativos às nove questões avaliadas por grupos de estudantes.
Nas barras dos gráficos, apresenta -se o percentual de alunos que assinalaram uma das opções ou a soma das porcentagens daqueles que assinalaram duas ( ou três ) delas.
Por exemplo, para as questões 1 e 2, nos gráficos, é apresentada a porcentagem total de participantes que assinalaram as opções ( D ) Difícil e ( E ) Muito difícil.
Em cada barra, foram assinalados, também, os extremos do intervalo de confiança de 95 \% como linhas verticais unidas por uma linha horizontal na forma da letra H maiúscula.
O estimador de um parâmetro com um certo nível de confiança ( por exemplo, 95,0 \% ) deve conter o parâmetro no intervalo de confiança em 95 \% das vezes.
Na comparação entre os estimadores dos parâmetros de duas classes de uma dada categoria ( por exemplo : Norte e Nordeste nas Grandes Regiões, de primeiro e último quarto dentro de desempenho, de IES Públicas e Privada, ou de Universidades e Faculdades ), associados aos seus respectivos intervalos de confiança, diz -se que não há diferença estatisticamente significativa entre os parâmetros das duas categorias quando há interseção entre os intervalos de confiança, e que há diferença, se os intervalos de confiança são disjuntos ( para maiores informações vide Glossário ).
Nas Tabelas do Anexo IIA ( Licenciatura ) e do Anexo IIB ( Bacharelado ) são apresentados os valores absolutos e a distribuição percentual18 das alternativas válidas das nove questões, segundo o mesmo recorte de desempenho dos alunos e da Grande Região de funcionamento do curso, de Categoria Administrativa e do tipo de Organização Acadêmica da IES. 4.1 GRAU DE DIFICULDADE DA PROVA 4.1.1 COMPONENTE DE FORMAÇÃO GERAL Ao avaliarem “ Qual o grau de dificuldade desta prova na parte de Formação Geral? ” ( Questão 1 ), 28,8 \% do grupo de inscritos e presentes de Licenciatura optaram pelas alternativas Difícil ou Muito difícil.
Entretanto, para mais da metade dos estudantes ( 61,1 \% ), o Componente de Formação Geral da prova foi avaliado como tendo grau de dificuldade Médio ( Gráfico 4.1a, Gráfico 4.2a e, no Anexo IIA, a Tabela II.1a ).
O percentual de estudantes de Licenciatura que consideraram a prova Difícil ou Muito difícil foi maior na região Norte, onde a proporção foi de 39,3 \%, enquanto a de menor proporção foi a Sudeste, com 24,7 \%.
No Gráfico 4.1a, é possível observar que as diferenças entre a região Norte e as demais regiões são estatisticamente significativas, assim como a diferença entre a região Sudeste e as demais regiões.
Nas Grandes Regiões, a proporção de presentes à prova que avaliaram o Componente de Formação Geral como de grau de dificuldade Médio esteve entre 55,0 \%, na região Norte, e 62,9 \%, na região Sudeste. entre a região Norte e as demais regiões são estatisticamente significativas, assim como a diferença entre a região Sudeste e as demais regiões.
Nas Grandes Regiões, a proporção de presentes à prova que avaliaram o Componente de Formação Geral como de grau de dificuldade Médio esteve entre 55,0 \%, na região Norte, e 62,9 \%, na região Sudeste.
Ao avaliarem “ Qual o grau de dificuldade desta prova na parte de Formação Geral? ” ( Questão 1 ), 17,3 \% do grupo de inscritos e presentes de Bacharelado optaram pelas alternativas Difícil ou Muito difícil.
Entretanto, entre mais da metade dos estudantes de Bacharelado ( 64,3 \% ), o Componente de Formação Geral da prova foi considerado como sendo de grau de dificuldade Médio ( Gráfico 4.1b, Gráfico 4.2b e, no Anexo IIB, a Tabela II.1b ).
O percentual de estudantes de Bacharelado que consideraram a prova como Difícil ou Muito difícil foi maior na região Norte, onde a proporção foi de 32,5 \%, enquanto a de menor incidência foi a Nordeste, com 14,7 \%.
No Gráfico 4.1b, é possível observar que as diferenças entre a região Norte e as demais regiões são estatisticamente significativas, assim como entre a região Sul e as regiões Nordeste e Sudeste.
Nas Grandes Regiões, a proporção de presentes à prova que consideraram o Componente de Formação Geral como sendo de grau de dificuldade Médio esteve entre 58,5 \%, na região Norte, e 68,3 \%, na região Nordeste.
O percentual de alunos de Licenciatura que consideraram o componente de Formação Geral da prova Difícil ou Muito difícil foi decrescente em função dos quartos de desempenho : 41,1 \%, no primeiro quarto, e 17,9 \%, no quarto, o grupo de melhor desempenho na prova.
Nos quartos de desempenho intermediários, a proporção de alunos que consideraram a prova Difícil ou Muito difícil foi de 32,5 \%, no segundo quarto, e de 24,0 \%, no terceiro quarto.
As diferenças entre os quartos de desempenho são estatisticamente significativas.
Para todos os quartos de desempenho, a alternativa modal para essa pergunta foi a Médio, com 52,0 \% e 66,1 \% dos respondentes nos quartos extremos, primeiro e quarto, respectivamente.
O percentual de alunos de Bacharelado que consideraram o componente de Formação Geral da prova Difícil ou Muito difícil foi decrescente em função dos quartos de desempenho : 23,3 \%, no primeiro quarto e 12,2 \%, no quarto, grupo de melhor desempenho na prova.
Nos quartos de desempenho intermediários, a proporção de alunos que consideraram a prova Difícil ou Muito difícil foi de 17,8 \%, no segundo quarto e de 15,9 \%, no terceiro quarto.
As diferenças entre o quarto inferior e os demais quartos de desempenho são estatisticamente significativas, assim como entre o quarto superior e os demais quartos de desempenho.
Em todos os quartos de desempenho, a alternativa modal para essa pergunta foi a Médio, com 65,9 \% e 61,3 \% dos respondentes nos quartos extremos, primeiro e quarto, respectivamente.
O percentual de estudantes de Licenciatura que consideraram o componente de Formação Geral da prova Difícil ou Muito difícil foi maior entre os estudantes de CEFET/IF, alcançando 37,9 \%, enquanto entre os estudantes das Universidades, a Organização Acadêmica com a menor incidência, alcançou 26,8 \%.
No Gráfico 4.3a, é possível observar que as diferenças entre os CEFTE/IF e as demais Organizações Acadêmicas, com exceção das Faculdades ( 32,1 \% ), são estatisticamente significativas.
Nas Organizações Acadêmicas, a proporção de presentes de Licenciatura à prova que avaliaram o Componente de Formação Geral como sendo de grau de dificuldade Médio esteve entre 55,2 \%, nos CEFET/IF, e 62,4 \%, nas Universidades ( ver também a Tabela II.2a, no Anexo IIA ).
O percentual de estudantes de Bacharelado que consideraram o componente de Formação Geral da prova como Difícil ou Muito difícil foi maior entre os estudantes de Centros Universitários, para os quais a proporção foi de 21,7 \%, enquanto as Universidades foram a Organização Acadêmica com a menor incidência, com 16,5 \%.
No Gráfico 4.3b, é possível observar que apenas a diferença entre Universidades e Centros Universitários é estatisticamente significativa.
Nas Organizações Acadêmicas, a proporção de presentes de Bacharelado à prova que consideraram o Componente de Formação Geral como sendo de grau de dificuldade Médio esteve entre 56,9 \%, nos CEFET/IF, e 68,8 \% nas Faculdades ( ver também a Tabela II.2b, no Anexo IIB ).
O percentual de estudantes de Licenciatura que consideraram o componente de Formação Geral da prova Difícil ou Muito difícil foi maior entre os estudantes de cursos de IES Privadas ( 29,9 \% ), sem diferença estatisticamente significativa da proporção registrada em IES Públicas ( 28,5 \% ).
Já a proporção de presentes de Licenciatura à prova que avaliaram o Componente de Formação Geral como de grau de dificuldade Médio foi de 61,3 \%, entre os estudantes de IES Públicas, e de 60,5 \%, entre os de IES Privadas ( ver Gráfico 4.4a ).
O percentual de estudantes de Bacharelado que consideraram o componente de Formação Geral da prova como Difícil ou Muito difícil foi maior entre os estudantes de cursos de IES Privadas ( 20,4 \% ), uma diferença estatisticamente significativa da proporção alcançada em IES Públicas ( 15,2 \% ).
Já a proporção de presentes de Bacharelado à prova que consideraram o Componente de Formação Geral como sendo de grau de dificuldade Médio foi de 63,2 \%, entre os estudantes de IES Públicas, e 65,9 \%, entre os de IES Privadas ( ver Gráfico 4.4b ). 4.1.2 COMPONENTE DE CONHECIMENTO ESPECÍFICO Ao responderem à Questão 2 – “ Qual o grau de dificuldade desta prova na parte de Componente Específico? ” – 36,9 \% do grupo de estudantes de Licenciatura classificaram -na como Difícil ou Muito difícil.
Além disso, o Componente modal de Conhecimento Específico da prova foi considerado como de grau de dificuldade Médio por 56,6 \% dos alunos ( Gráfico 4.5a, Gráfico 4.6a, e, no Anexo IIA, a Tabela II.3a ).
A análise das respostas dos estudantes de Licenciatura quanto ao grau de dificuldade do Componente de Conhecimento Específico da prova, agregado por Grande Região, mostra diferenças estatisticamente significativas entre as proporções de alunos que a avaliaram como Difícil ou Muito difícil.
Tais diferenças são observadas entre a região com a maior proporção, a Norte ( 44,1 \% ), e as regiões Nordeste ( 36,7 \% ), Sudeste ( 35,1 \% ), a menor proporção, e Centro-Oeste ( 35,4 \% ).
A diferença entre as regiões Sudeste e Sul ( 39,8 \% ) também é estatisticamente significativa.
O percentual de alunos que classificaram o grau de dificuldade como Médio, no Componente de Conhecimento Específico, variou de 50,5 \% a 57,9 \%, nas regiões Norte, a menor proporção, e nas regiões Nordeste e Centro-Oeste, as maiores proporções, respectivamente.
Ao responderem à Questão 2 – “ Qual o grau de dificuldade desta prova na parte de Componente Específico? ” – 33,9 \% do grupo de estudantes de Bacharelado classificaram -na como Difícil ou Muito difícil.
Além disso, o Componente modal de Conhecimento Específico da prova foi considerado como sendo de grau de dificuldade Médio por 60,0 \% dos alunos ( Gráfico 4.5b, Gráfico 4.6b, e, no Anexo IIB, a Tabela II.3b ).
A análise das respostas dos estudantes de Bacharelado quanto ao grau de dificuldade do Componente de Conhecimento Específico da prova, agregado por Grande Região, mostra que a diferença entre a maior, na região Norte ( 45,8 \% ), e todas as outras proporções de alunos que a avaliaram como Difícil ou Muito difícil, é estatisticamente significativa, exceto na região Sul ( 37,6 \% ).
As diferenças entre a região Nordeste ( 31,6 \% ), a menor proporção, e a região Sul, também são estatisticamente significativas, assim como a diferença entre as regiões Sudeste ( 32,9 \% ) e Sul.
O percentual de alunos que classificaram o grau de dificuldade como Médio, no Componente de Conhecimento Específico, variou de 50,0 \% a 61,3 \%, nas regiões Norte e Nordeste, respectivamente.
Componente de Conhecimento Específico, variou de 50,0 \% a 61,3 \%, nas regiões Norte e Nordeste, respectivamente.
Considerando -se a avaliação da dificuldade das questões do Componente de Conhecimento Específico da prova de acordo com o desempenho dos estudantes de Licenciatura, observa -se uma tendência decrescente dos que consideraram a parte específica Difícil ou Muito difícil, variando de 28,7 \% ( último quarto ) a 43,8 \% ( primeiro quarto ).
As demais proporções foram de 40,7 \% e 34,6 \%, respectivamente, no segundo e terceiro quartos.
Observa -se que há diferença estatisticamente significativa dos resultados entre o quarto inferior de desempenho e os dois quartos superiores de desempenho, assim como entre o segundo quarto e os dois quartos superiores de desempenho, além da diferença entre os dois quartos superiores.
Já os que responderam que o grau de dificuldade das questões do Componente de Conhecimento Específico da prova foi Médio atingiram percentagens que variaram de 51,1 \%, no primeiro quarto, a 62,0 \%, no último.
Considerando -se a avaliação da dificuldade das questões do Componente de Conhecimento Específico da prova, de acordo com o desempenho dos estudantes de Bacharelado, observa -se que só há diferença estatisticamente significativa dos resultados entre o quarto superior de desempenho e os demais quartos, assim como a diferença entre o primeiro e terceiro quartos.
A proporção dos que classificaram a parte específica como Difícil ou Muito difícil variou de 28,5 \% ( último quarto ) a 38,0 \% ( primeiro quarto ).
As demais proporções dessa reposta foram de 35,8 \% e 33,2 \%, nos quartos segundo e terceiro, respectivamente.
Já a proporção dos que responderam que o grau de dificuldade das questões do Componente de Conhecimento Específico da prova foi Médio variou de 56,8 \%, no primeiro quarto, a 63,6 \%, no último.
Considerando -se a avaliação da dificuldade das questões do Componente de Conhecimento Específico da prova, o percentual de estudantes de Licenciatura que as consideraram Difícil ou Muito difícil foi maior entre os estudantes de CEFET/IF, atingindo 44,7 \% ; já entre os de Universidades, a Organização Acadêmica com a menor incidência, essa proporção foi de 35,1 \%.
No Gráfico 4.7a, é possível observar que a diferença entre CEFET/IF e as demais Organizações Acadêmicas, com exceção das Faculdades ( 40,6 \% ), é estatisticamente significativa.
Nas Organizações Acadêmicas, a proporção de presentes de Licenciatura à prova que consideraram o Componente de Conhecimento Específico de grau de dificuldade Médio esteve entre 50,3 \%, nos CEFET/IF, e 58,0 \%, nas Universidades ( ver também Tabela II.4a, no Anexo IIA ).
Considerando -se a avaliação da dificuldade das questões do Componente de Conhecimento Específico da prova, o percentual de estudantes de Bacharelado que as consideraram como Difícil ou Muito difícil foi maior entre os estudantes de Faculdades, entre os quais a proporção foi de 39,0 \% ; enquanto os CEFET/IF foram a Organização Acadêmica com a menor incidência, com a proporção de 31,9 \%.
Existe diferença estatisticamente significativa apenas entre as Universidades ( 33,0 \% ) e os Centros Universitários ( 38,3 \% ).
Nas Organizações Acadêmicas, a proporção de presentes de Bacharelado à prova que consideraram o Componente de Formação Geral como sendo de grau de dificuldade Médio esteve entre 56,5 \%, nos Centros Universitários, e 66,7 \% nos CEFET/IF ( ver também Tabela II.4b, no Anexo IIB ).
O percentual de estudantes de Licenciatura que consideraram o Componente Específico da prova Difícil ou Muito difícil foi maior entre os estudantes de cursos de IES Privadas ( 40,1 \% ), com diferença estatisticamente significativa da proporção registrada em IES Públicas ( 35,9 \% ).
Já a proporção de presentes de Licenciatura à prova que avaliaram esse componente como sendo de grau de dificuldade Médio foi de 57,4 \%, entre os estudantes de IES Públicas, e de 53,9 \%, entre os de IES Privadas ( ver Gráfico 4.8a e Tabela II.4a, no Anexo IIA ).
O percentual de estudantes de Bacharelado que consideraram o Componente Específico da prova como Difícil ou Muito difícil foi maior entre os estudantes de cursos de IES Privadas ( 38,1 \% ), uma diferença estatisticamente significativa da proporção em IES Públicas ( 31,0 \% ).
Já a proporção de presentes de Bacharelado à prova que consideraram esse componente como sendo de grau de dificuldade Médio foi de 61,9 \%, entre os estudantes de IES Públicas, e de 57,3 \%, entre os de IES Privadas ( ver Gráfico 4.8b e Tabela II.4b, no Anexo IIB ). 4.2 EXTENSÃO DA PROVA EM RELAÇÃO AO TEMPO TOTAL Indagados quanto à extensão da prova, em relação ao tempo total oferecido para a sua resolução ( Questão 3 ), os estudantes de Licenciatura apontaram, com maior incidência, a alternativa que considerava a extensão Adequada, em todas as agregações consideradas ( Gráfico 4.9a, Gráfico 4.10a, e, no Anexo IIA, a Tabela II.5a ).
O percentual de alunos de Licenciatura que responderam ser a extensão da prova Adequada foi de 56,2 \%.
Já 36,7 \% dos inscritos presentes consideraram que a prova foi Longa ou Muito longa, e 7,1 \% a avaliaram como Curta ou Muito curta.
Entre as Grandes Regiões, os que avaliaram a prova como Longa ou Muito longa, em relação ao tempo total destinado à sua resolução, atingiram uma porcentagem que variou de 31,4 \%, na região Norte, até 39,2 \%, na região Sul.
Há diferença estatisticamente significativa das regiões Sudeste e Sul em relação às demais regiões, mas não entre si.
Indagados quanto à extensão da prova, em relação ao tempo total oferecido para a sua resolução ( Questão 3 ), os estudantes de Bacharelado apontaram, com maior incidência, a alternativa que considerava a extensão Adequada, em todas as agregações consideradas ( Gráfico 4.9b, Gráfico 4.10b, e, no Anexo IIB, a Tabela II.5b ).
O percentual de alunos de Bacharelado que responderam ser a extensão da prova Adequada foi de 60,7 \%.
Já 31,9 \% dos inscritos presentes consideraram que a prova foi Longa ou Muito longa, e 7,4 \% a avaliaram como Curta ou Muito curta.
Entre as Grandes Regiões, a proporção daqueles que avaliaram a prova como Longa ou Muito longa, em relação ao tempo total destinado à sua resolução, variou de 26,2 \%, na região Centro-Oeste, até 34,1 \%, na região Nordeste.
Há diferença estatisticamente significativa entre a região Centro-Oeste e as regiões Nordeste, Sudeste e Sul.
Considerando -se o desempenho dos alunos de Licenciatura, nota -se ainda que 53,3 \% consideraram a extensão da prova Adequada, no quarto de desempenho inferior, e 58,7 \%, no de melhor desempenho ( quarto superior ).
Nos quartos intermediários, essa proporção foi de 55,4 \%, no segundo quarto, e de 57,2 \%, no terceiro.
No Gráfico 4.10a, observa -se que as proporções de estudantes de Licenciatura que consideraram a prova Longa ou Muito longa apresentam tendência decrescente em função dos quartos de desempenho, sendo a maior proporção dessa resposta no primeiro quarto de desempenho ( 38,0 \% ), e a menor, no quarto superior ( 34,8 \% ).
Pode -se constatar que há diferença estatisticamente significativa apenas entre a proporção do quarto inferior de desempenho e a do quarto superior.
Considerando -se o desempenho dos alunos de Bacharelado, nota -se ainda que 58,8 \% consideraram a extensão da prova Adequada, no quarto de desempenho inferior, e 60,9 \%, no de melhor desempenho ( quarto superior ).
Nos quartos intermediários, essa proporção foi de 62,2 \%, no segundo, e de 60,7 \%, no terceiro quarto.
No Gráfico 4.10b, pode -se constatar que não há diferença estatisticamente significativa entre as proporções de estudantes de Bacharelado que consideraram a prova Longa ou Muito longa em função dos quartos de desempenho, sendo a maior proporção dessa resposta no primeiro quarto de desempenho ( 33,6 \% ), e a menor, no segundo quarto ( 30,2 \% ).
Considerando -se a extensão da prova, o percentual de estudantes de Licenciatura que avaliaram a prova como Longa ou Muito longa, em relação ao tempo total destinado à sua resolução, foi maior entre os estudantes de Centros Universitários, alcançando 39,0 \%, enquanto entre os estudantes de Faculdades, a Organização Acadêmica com a menor incidência, esse índice atingiu 32,7 \%.
No Gráfico 4.11a, é possível observar que não há diferença estatisticamente significativa entre os tipos de Organização Acadêmica.
A proporção de presentes de Licenciatura à prova que consideraram a sua extensão Adequada, em relação ao tempo total destinado à sua resolução, esteve entre 54,6 \%, nos Centros Universitários, e 59,2 \%, nas Faculdades ( ver também Tabela II.6a, no Anexo IIA ).
Considerando -se a extensão da prova, o percentual de estudantes de Bacharelado que avaliaram a prova como Longa ou Muito longa, em relação ao tempo total destinado à sua resolução foi maior entre os estudantes de CEFET/IF, para os quais a proporção foi de 36,1 \%, enquanto os de Centros Universitários foram a Organização Acadêmica com a menor incidência, com 30,4 \%.
No Gráfico 4.11b, é possível observar que não há diferenças estatisticamente significativas entre os tipos de Organização Acadêmica.
Nas Organizações Acadêmicas, a proporção de presentes de Bacharelado à prova que consideraram a sua extensão como Adequada, em relação ao tempo total destinado à sua resolução, esteve entre 55,6 \%, nos CEFET/IF, e 63,3 \%, nas Faculdades ( ver também Tabela II.6b, no Anexo IIB ).
O percentual de estudantes de Licenciatura que consideraram a prova Longa ou Muito longa, em relação ao tempo total destinado à sua resolução, foi maior entre os estudantes de cursos de IES Privadas ( 38,3 \% ), não havendo diferença estatisticamente significativa da proporção em relação aos de IES Públicas ( 36,2 \% ).
Já a proporção de presentes de Licenciatura à prova que consideraram a extensão Adequada foi de 56,6 \% entre os estudantes de IES Públicas, e de 55,0 \%, entre os de IES Privadas ( ver Gráfico 4.12a e Tabela II.6a, no Anexo IIA ).
O percentual de estudantes de Bacharelado que consideraram a prova como Longa ou Muito longa, em relação ao tempo total destinado à sua resolução foi maior entre os estudantes de cursos de IES Públicas ( 33,2 \% ), com diferença estatisticamente significativa da proporção registrada em IES Privadas ( 29,9 \% ).
Já a proporção de presentes de Bacharelado à prova que consideraram a extensão como Adequada foi de 59,8 \%, entre os de IES Públicas, e de 62,0 \%, entre os de IES Privadas ( ver Gráfico 4.12b e Tabela II.6b, no Anexo IIB ). 4.3 COMPREENSÃO DOS ENUNCIADOS DAS QUESTÕES 4.3.1 COMPONENTE DE FORMAÇÃO GERAL Em relação aos enunciados das questões do Componente de Formação Geral ( Questão 4 ), as opiniões foram positivas, já que 77,3 \% dos alunos de Licenciatura avaliados consideraram todos ou a maioria dos enunciados das questões “ claros e objetivos ” ( Gráfico 4.13a, Gráfico 4.14a e, no Anexo IIA, a Tabela II.7a ).
Na análise regional, a percentagem de estudantes de Licenciatura que avaliaram que todos ou a maioria dos enunciados das questões do Componente de Formação Geral estavam “ claros e objetivos ” variou de 70,8 \%, na região Norte, a 80,9 \%, na região Sul.
A análise das percepções dos estudantes de Licenciatura sobre a clareza e a objetividade dos enunciados permite afirmar que todos ou a maioria dos enunciados de questões relativas ao Componente de Formação Geral, foram considerados “ claros e objetivos ” para a maior parte dos respondentes ( maior ou igual a 70,8 \%, em todas as regiões, e maior ou igual a 67,5 \%, em todos os quartos de desempenho ).
Em relação aos enunciados das questões do Componente de Formação Geral ( Questão 4 ), as opiniões foram positivas, já que 80,7 \% dos alunos de Bacharelado avaliados consideraram todos ou a maioria dos enunciados das questões “ claros e objetivos ” ( Gráfico 4.13b, Gráfico 4.14b, e, no Anexo IIB, a Tabela II.7b ).
Na análise regional, a percentagem de estudantes de Bacharelado que avaliaram que todos ou a maioria dos enunciados das questões do Componente de Formação Geral estavam “ claros e objetivos ” variou de 76,9 \%, na região Norte, a 82,8 \%, na região Centro-Oeste.
A análise das percepções dos estudantes de Bacharelado sobre a clareza e objetividade dos enunciados permite afirmar que todos ou a maioria dos enunciados de questões relativas ao Componente de Formação Geral, foram considerados “ claros e objetivos ” entre a maior parte dos respondentes ( maior ou igual a 76,9 \%, em todas as regiões, e maior ou igual a 75,8 \%, em todos os quartos de desempenho ).
Segundo o desempenho, observa -se que a proporção dos que emitiram essa opinião cresce, conforme o desempenho aumenta, com diferenças estatisticamente significativas entre os quartos de desempenho.
No quarto superior, a clareza e a objetividade de todos ou da maioria dos enunciados das questões foram percebidas por 84,8 \%.
Segundo o desempenho, observa -se que a proporção dos que emitiram essa opinião cresce, conforme o desempenho aumenta, até o terceiro quarto, com diferenças estatisticamente significativas apenas entre o quarto inferior e os demais quartos de desempenho.
No quarto superior, a clareza e a objetividade de todos ou da maioria dos enunciados das questões foram percebidas por 82,2 \%.
A proporção dos presentes de Licenciatura que consideraram que todos ou a maioria dos enunciados das questões do Componente de Formação Geral estavam “ claros e objetivos ” foi maior entre os estudantes de Centros Universitários, chegando a 80,1 \%, enquanto entre os de CEFET/IF, a Organização Acadêmica com a menor incidência, esse índice atingiu 73,4 \%.
No Gráfico 4.15a, é possível observar que as diferenças entre os CEFET/IF e os demais tipos de Organização Acadêmica, com exceção das Faculdades ( 78,9 \% ), são estatisticamente significativas ( ver também Tabela II.8a, no Anexo IIA ).
A proporção dos presentes de Bacharelado que consideraram que todos ou a maioria dos enunciados das questões do Componente de Formação Geral estavam “ claros e objetivos ” foi maior entre os estudantes de Faculdades, entre os quais a proporção foi de 86,1 \%, enquanto entre os de CEFET/IF, a Organização Acadêmica com a menor incidência, o índice alcançou 72,2 \%.
No Gráfico 4.15b, é possível observar que há diferença estatisticamente significativa apenas entre as Universidades ( 80,5 \% ) e as Faculdades ( ver também Tabela II.8b, no Anexo IIB ).
O percentual de estudantes de Licenciatura que consideraram que todos ou a maioria dos enunciados das questões do Componente de Formação Geral estavam “ claros e objetivos ” foi maior entre os estudantes de cursos de IES Privadas ( 79,1 \% ), com diferença estatisticamente significativa da proporção registrada em IES Públicas ( 76,8 \% ).
Um maior detalhamento pode ser encontrado no Gráfico 4.16a e na Tabela II.8a, no Anexo IIA.
O percentual de estudantes de Bacharelado que consideraram que todos ou a maioria dos enunciados das questões do Componente de Formação Geral estavam “ claros e objetivos ” foi maior entre os estudantes de cursos de IES Privadas ( 82,1 \% ), sem diferença estatisticamente significativa da proporção em IES Públicas ( 79,8 \% ).
Um maior detalhamento pode ser encontrado no Gráfico 4.16b e na Tabela II.8b, no Anexo IIB. 4.3.2 COMPONENTE DE CONHECIMENTO ESPECÍFICO Em relação aos enunciados das questões do Componente de Conhecimento Específico da prova, para 79,6 \% dos estudantes de Licenciatura avaliados da Área de Ciências Biológicas, a clareza e a objetividade ( Questão 5 ) estavam presentes em todas ou na maioria das questões ( Gráfico 4.17a a Gráfico 4.20a, e no Anexo IIA, as Tabelas II.9a e II.10a ).
A maioria dos estudantes de Licenciatura de todas as Grandes Regiões brasileiras considerou “ claros e objetivos ” todos ou a maioria dos enunciados das questões do Componente de Conhecimento Específico da prova, percentual sempre maior ou igual a 72,9 \%.
A diferença entre a região Norte ( 72,9 \% ) e as demais regiões é estatisticamente significativa, bem como entre a região Nordeste ( 77,4 \% ) e as demais regiões.
Em relação aos enunciados das questões do Componente de Conhecimento Específico da prova, para 83,6 \% dos estudantes de Bacharelado avaliados da Área de Ciências Biológicas, a clareza e a objetividade ( Questão 5 ) estavam presentes em todas ou na maioria das questões ( Gráfico 4.17b a Gráfico 4.20b, e no Anexo IIB, as Tabelas II.9b e II.10b ).
A maioria dos estudantes de Bacharelado de todas as Grandes Regiões brasileiras considerou “ claros e objetivos ” todos ou a maioria dos enunciados das questões do Componente de Conhecimento Específico da prova, percentual sempre maior ou igual a 81,0 \%.
Não há diferença estatisticamente significativa entre as regiões.
A proporção de estudantes de Licenciatura que consideraram os enunciados das questões “ claros e objetivos ” apresenta uma tendência crescente em relação ao aumento de desempenho : mais elevada no quarto superior ( 87,7 \% ), se comparada ao quarto inferior de desempenho ( 69,2 \% ).
As diferenças entre todos os quartos de desempenho são estatisticamente significativas.
A proporção de estudantes de Bacharelado que consideraram os enunciados das questões “ claros e objetivos ” apresenta uma tendência crescente em relação ao aumento de desempenho : mais elevada no quarto superior ( 86,3 \% ), se comparada ao quarto inferior de desempenho ( 78,5 \% ).
As diferenças entre o quarto inferior e os demais quartos de desempenho são estatisticamente significativas.
A proporção dos presentes de Licenciatura que consideraram que todos ou a maioria dos enunciados das questões do Componente de Conhecimento Específico estavam “ claros e objetivos ” foi maior entre os estudantes de Centros Universitários, atingindo 80,5 \%.
Entre os estudantes de Faculdades, a Organização Acadêmica com a menor incidência, essa proporção atingiu 76,0 \%.
No Gráfico 4.19a, é possível observar que apenas a diferença entre o CEFET/IF ( 76,3 \% ) e as Universidades ( 80,4 \% ) é estatisticamente significativa ( ver também Tabela II.10a, no Anexo IIA ).
A proporção dos presentes de Bacharelado que consideraram que todos ou a maioria dos enunciados das questões do Componente de Conhecimento Específico estavam “ claros e objetivos ” foi maior entre os estudantes de CEFET/IF, para os quais a proporção foi de 87,5 \%, enquanto entre os de Centros Universitários, a Organização Acadêmica com a menor incidência, o índice atingiu 83,3 \%.
No Gráfico 4.19b, é possível observar que não há diferença estatisticamente significativa entre os tipos de Organizações Acadêmicas ( ver também Tabela II.10b, no Anexo IIB ).
O percentual de estudantes de Licenciatura que consideraram que todos ou a maioria dos enunciados das questões do Componente de Conhecimento Específico estavam “ claros e objetivos ” foi maior entre os estudantes de cursos de IES Privadas ( 79,7 \% ), sem diferença estatisticamente significativa da proporção registrada em IES Públicas ( 79,6 \% ).
Ver Gráfico 4.20a e Tabela II.10a, no Anexo IIA, para maior detalhamento.
O percentual de estudantes de Bacharelado que consideraram que todos ou a maioria dos enunciados das questões do Componente de Conhecimento Específico estavam “ claros e objetivos ” foi menor entre os estudantes de cursos de IES Privadas ( 83,2 \% ), sem diferença estatisticamente significativa da proporção registrada em IES Públicas ( 83,8 \% ).
Ver Gráfico 4.20b e Tabela II.10b, no Anexo IIB, para um maior detalhamento. 4.4 SUFICIÊNCIA DAS INFORMAÇÕES/INSTRUÇÕES FORNECIDAS Ao se avaliarem as informações/instruções fornecidas para a resolução das questões ( Questão 6 ), constatou -se que 84,0 \% dos respondentes de Licenciatura da Área de Ciências Biológicas de todo o Brasil afirmaram que essas eram até excessivas ou suficientes em todas ou na maioria das questões ( Gráfico 4.21a, Gráfico 4.22a e, no Anexo IIA, a Tabela II.11a ).
Quanto à distribuição de respondentes pelas Grandes Regiões, observa -se que a proporção de estudantes que consideraram as informações/instruções fornecidas até excessivas ou suficientes em todas ou na maioria das questões foi sempre superior ou igual a 76,0 \%, na região Norte, chegando a 86,7 \%, na região Sul.
Existem diferenças estatisticamente significativas entre a região Norte e as demais regiões e entre a região Nordeste ( 81,8 \% ) e as demais regiões.
Ao avaliarem as informações/instruções fornecidas para a resolução das questões ( Questão 6 ), 87,8 \% dos respondentes de Bacharelado da Área de Ciências Biológicas de todo o Brasil afirmaram que essas eram até excessivas ou suficientes em todas ou na maioria das questões ( Gráfico 4.21b, Gráfico 4.22b, e, no Anexo IIB, a Tabela II.11b ).
Quanto à distribuição de respondentes pelas Grandes Regiões, observa -se que a proporção de estudantes que consideraram as informações/instruções fornecidas até excessivas ou suficientes em todas ou na maioria das questões foi sempre superior ou igual a 81,5 \%, na região Norte, chegando a 90,3 \%, na região Centro-Oeste.
Existem diferenças estatisticamente significativas apenas entre a região Norte e as regiões Sul e Centro-Oeste.
Levando -se em conta o desempenho dos participantes de Licenciatura, observa -se uma tendência crescente das proporções de participantes que avaliaram as informações/instruções como até excessivas ou suficientes em todas ou na maioria das questões, havendo diferenças estatisticamente significativas entre os quartos de desempenho, conforme mostrado no Gráfico 4.22a.
O percentual foi mais elevado no quarto superior ( 90,7 \% ), percentual superior ao da média nacional ( 84,0 \% ).
Já no quarto inferior, a suficiência das informações/instruções declarada como até excessiva, em todas ou na maioria das questões foi percebida por 75,3 \% dos respondentes.
Levando -se em conta o desempenho dos participantes de Bacharelado, observa -se uma tendência crescente das proporções de participantes que avaliaram as informações/instruções como até excessivas ou suficientes em todas ou na maioria das questões, havendo diferenças estatisticamente significativas entre o quarto inferior e os demais quartos de desempenho, e entre o segundo quarto e o quarto superior de desempenho, como mostrado no Gráfico 4.22b.
O percentual foi mais elevado no quarto superior ( 90,8 \% ), percentual superior à média nacional ( 87,8 \% ).
Já no quarto inferior, a suficiência das informações/instruções declarada como até excessiva, em todas ou na maioria das questões foi percebida por 83,9 \% dos respondentes.
Considerando -se as informações/instruções fornecidas para a resolução das questões ( Questão 6 ), a proporção que considerou que eram até excessivas ou suficientes em todas ou na maioria das questões foi maior entre os estudantes de Licenciatura de Centros Universitários, totalizando 87,7 \%.
Os CEFET/IF foram a Organização Acadêmica em que houve a menor proporção : 79,5 \%.
No Gráfico 4.23a, é possível observar que as diferenças entre as Universidades ( 84,7 \% ), em relação aos Centros Universitários e aos CEFET/IF, são estatisticamente significativas, assim como entre os Centros Universitários e os CEFET/IF ( ver também Tabela II.12a, no Anexo IIA ).
Considerando -se as informações/instruções fornecidas para a resolução das questões ( Questão 6 ), a proporção que considerou que eram até excessivas ou suficientes em todas ou na maioria das questões foi maior entre os estudantes de Bacharelado de CEFET/IF, para os quais a proporção foi de 88,9 \%, enquanto entre os de Centros Universitários, a Organização Acadêmica com a menor incidência, o índice atingiu 86,8 \%.
No Gráfico 4.23b, é possível observar que não existe diferença estatisticamente significativa entre os tipos de Organização Acadêmica ( ver também Tabela II.12b, no Anexo IIB ).
O percentual de estudantes de Licenciatura que consideraram que as informações/instruções fornecidas para a resolução das questões eram até excessivas ou suficientes em todas ou na maioria das questões foi menor entre os estudantes de cursos de IES Públicas ( 83,8 \% ), não havendo diferença estatisticamente significativa da proporção registrada em IES Privadas ( 84,8 \% ) ( ver Gráfico 4.24a e Tabela II.12a, no Anexo IIA ).
O percentual de estudantes de Bacharelado que consideraram que as informações/instruções fornecidas para a resolução das questões eram até excessivas ou suficientes em todas ou na maioria das questões foi maior entre os estudantes de cursos de IES Públicas ( 88,2 \% ), sem diferença estatisticamente significativa da proporção registrada em IES Privadas ( 87,2 \% ) ( ver Gráfico 4.24b e Tabela II.12b, no Anexo IIB ). 4.5 DIFICULDADE ENCONTRADA AO RESPONDER À PROVA Indagados sobre as dificuldades com as quais se depararam ao responder à prova ( Questão 7 ), 20,0 \% dos estudantes de Licenciatura apontaram o Desconhecimento do conteúdo.
Para 44,8 \%, a Forma diferente de abordagem do conteúdo foi indicada como dificuldade.
Já a Falta de motivação para fazer a prova foi a dificuldade apontada por 17,9 \% dos respondentes.
Considerando -se todo o Brasil, 13,5 \% dos respondentes afirmaram que não tiveram qualquer tipo de dificuldade para responder à prova ( Tabelas II.13a e II.14a, no Anexo IIA ).
Nos Gráficos 4.25a a 4.28a, são apresentados os percentuais de estudantes de Licenciatura que apontaram o Desconhecimento do conteúdo como “ dificuldade ao responder à prova ”.
Na análise por Grandes Regiões, o percentual de inscritos e presentes de Licenciatura que apontaram o Desconhecimento do conteúdo como “ dificuldade ao responder à prova ” não superou 22,8 \%.
Os percentuais variaram de 19,2 \%, na região Nordeste, a 22,8 \%, na região Sul.
A Forma diferente de abordagem do conteúdo foi a escolha modal dos estudantes de Licenciatura, com percentuais que variaram de 40,7 \% ( região Sudeste ) a 55,3 \% ( região Norte ).
O percentual de alunos que citaram a Falta de motivação como dificuldade variou de 12,0 \% ( região Norte ) a 19,6 \% ( região Sudeste ).
Os que declararam não ter qualquer tipo de dificuldade para responder à prova variaram de 9,4 \%, na região Norte, a 15,6 \%, na região Sudeste.
Indagados sobre as dificuldades com as quais se depararam ao responder à prova ( Questão 7 ), 22,3 \% dos estudantes de Bacharelado apontaram o Desconhecimento do conteúdo.
Para 41,2 \%, a Forma diferente de abordagem do conteúdo foi indicada como dificuldade.
Já a Falta de motivação para fazer a prova foi a dificuldade apontada por 20,8 \% dos respondentes.
Considerando -se todo o Brasil, 12,9 \% dos respondentes afirmaram que não tiveram qualquer tipo de dificuldade para responder à prova ( Tabelas II.13b e II.14b, no Anexo IIB ).
Nos Gráficos 4.25b a 4.28b, são apresentados os percentuais de estudantes de Bacharelado que apontaram o Desconhecimento do conteúdo como “ dificuldade ao responder à prova ”.
Na análise por Grandes Regiões, o percentual de inscritos e presentes de Bacharelado que apontaram o Desconhecimento do conteúdo como “ dificuldade ao responder à prova ” não superou 24,4 \%.
Os percentuais variaram de 15,2 \%, na região Norte, a 24,4 \%, na região Sul.
A Forma diferente de abordagem do conteúdo foi a escolha modal dos estudantes de Bacharelado, com percentuais que variaram de 39,8 \% ( região Sul ) a 60,7 \% ( região Norte ).
O percentual de alunos que citaram a Falta de motivação como dificuldade variou de 13,3 \% ( região Norte ) a 27,8 \% ( região Nordeste ).
Os que declararam não ter qualquer tipo de dificuldade para responder à prova variaram de 10,0 \%, na região Norte, a 14,2 \%, na região Sudeste.
Em relação aos quartos de desempenho, o Desconhecimento do conteúdo foi a opção escolhida por 17,8 \% dos estudantes de Licenciatura do quarto inferior e por 23,0 \% do quarto superior.
São observadas diferenças estatisticamente significativas entre o quarto superior e os demais quartos de desempenho.
A alternativa modal entre os alunos de Licenciatura, quando agregados pelos quartos de desempenho, para a dificuldade encontrada, foi causada pela Forma diferente de abordagem do conteúdo : 48,5 \%, no primeiro quarto, e 40,1 \% do quarto superior, escolheram essa alternativa.
Em relação aos quartos de desempenho, o Desconhecimento do conteúdo foi a opção escolhida por 18,9 \% dos estudantes de Bacharelado do quarto inferior e por 25,4 \% do quarto superior.
É observada diferença estatisticamente significativa apenas entre o quarto inferior e o segundo e quarto quartos de desempenho.
A alternativa modal entre os alunos de Bacharelado, quando agregados pelos quartos de desempenho, para a dificuldade encontrada, foi causada pela Forma diferente de abordagem do conteúdo : variando de 39,2 \%, no quarto superior, a 42,9 \%, no primeiro quarto.
Na análise por Organização Acadêmica, o percentual de inscritos e presentes de Licenciatura que apontaram o Desconhecimento do conteúdo como “ dificuldade ao responder à prova ” não superou 22,5 \%.
Os percentuais variaram de 19,0 \%, nas Faculdades, a 22,5 \%, nos Centros Universitários.
A Forma diferente de abordagem do conteúdo foi a escolha modal dos estudantes de Licenciatura, com percentuais que variaram de 41,7 \% ( Centros Universitários ) a 50,7 \% ( CEFET/IF ).
O percentual de alunos de Licenciatura que citaram a Falta de motivação como dificuldade variou de 12,0 \% ( Faculdades ) a 19,0 \% ( Universidades ).
Os que declararam não ter qualquer tipo de dificuldade para responder à prova variaram de 12,2 \%, nos CEFET/IF, a 17,6 \%, nos Centros Universitários.
No Gráfico 4.27a, é possível observar que não há diferenças estatisticamente significativas entre os tipos de Organização Acadêmica ( ver também Tabela II.14a, no Anexo IIA ).
Na análise por Organização Acadêmica, o percentual de inscritos e presentes de Bacharelado que apontaram o Desconhecimento do conteúdo como “ dificuldade ao responder à prova ” não superou 25,0 \%.
Os percentuais variaram de 20,0 \%, nas Faculdades, a 25,0 \%, nos CEFET/IF.
A Forma diferente de abordagem do conteúdo foi a escolha modal dos estudantes de Bacharelado, com percentuais que variaram de 36,1 \% ( CEFET/IF ) a 44,4 \% ( Centros Universitários ).
O percentual de alunos de Bacharelado que citaram a Falta de motivação como dificuldade variou de 14,4 \% ( Centros Universitários ) a 22,0 \% ( Universidades ).
Os que declararam não ter qualquer tipo de dificuldade para responder à prova variaram de 12,2 \%, nas Universidades a 19,6 \%, nas Faculdades.
No Gráfico 4.27b, é possível observar que não há diferenças estatisticamente significativas entre os tipos de Organização Acadêmica ( ver também Tabela II.14b, no Anexo IIB ).
Em relação às Categorias Administrativas, o Desconhecimento do conteúdo foi a opção escolhida por 19,5 \% dos estudantes de Licenciatura, nas IES Públicas, e por 21,6 \%, nas IES Privadas.
Não são observadas diferenças estatisticamente significativas entre essas duas categorias.
A alternativa modal entre os alunos de Licenciatura, quando agregados pelas Categorias Administrativas, para a dificuldade encontrada, foi causada pela Forma diferente de abordagem do conteúdo : 45,2 \%, entre os de IES Públicas, e 43,6 \%, entre os IES Privadas ( ver Gráfico 4.28a e Tabela II.14a, no Anexo IIA ).
Em relação às Categorias Administrativas, o Desconhecimento do conteúdo foi a opção escolhida por 21,3 \% dos estudantes de Bacharelado, nas IES Públicas, e por 23,8 \%, nas IES Privadas.
Não são observadas diferenças estatisticamente significativas entre essas duas categorias.
A alternativa modal entre os alunos de Bacharelado, quando agregados pelas Categorias Administrativas, para a dificuldade encontrada, foi causada pela Forma diferente de abordagem do conteúdo : 39,8 \%, entre os de IES Públicas, e 43,4 \%, entre os de IES Privadas ( ver Gráfico 4.28b e Tabela II.14b, no Anexo IIB ). 4.6 CONTEÚDOS DAS QUESTÕES OBJETIVAS DA PROVA Ao se analisarem os conteúdos das questões objetivas da prova ( Questão 8 ), verifica -se que um percentual muito pequeno dos estudantes de Licenciatura avaliados, apenas 4,3 \%, afirmou que não estudou ainda a maioria desses conteúdos ( Gráficos 4.29a a Gráfico 4.32a, e as Tabelas II.15a e II.16a, no Anexo IIA ).
A maioria ( 66,2 \% ) afirmou ter estudado e aprendido muitos ou todos os conteúdos avaliados.
Na análise por Grande Região, a proporção de respondentes de Licenciatura que escolheram a opção não estudou ainda a maioria desses conteúdos foi pequena.
Observa -se que, nas regiões Norte ( 6,8 \% ) e Nordeste ( 4,4 \% ), as proporções foram maiores do que a da média nacional ( 4,3 \% ).
Nas regiões Sudeste, Sul e Centro-Oeste, a situação foi inversa, com proporções mais baixas, correspondendo a, respectivamente, 3,9 \%, 4,1 \% e 3,4 \%.
São observadas diferenças estatisticamente significativas apenas entre a região Norte e as demais regiões.
Em todas as regiões, a maioria dos presentes afirmou ter estudado e aprendido muitos ou todos os conteúdos, com proporções variando entre 50,9 \%, na região Norte, a 70,2 \%, na região Sudeste.
Ao analisar os conteúdos das questões objetivas da prova ( Questão 8 ), um percentual muito pequeno dos estudantes de Bacharelado avaliados, apenas 2,2 \%, afirmaram que não estudou ainda a maioria desses conteúdos ( Gráficos 4.29b a Gráfico 4.32b, e as Tabelas II.15b e II.16b, no Anexo IIB ).
A maioria ( 72,6 \% ) afirmou ter estudado e aprendido muitos ou todos os conteúdos avaliados.
Na análise por Grande Região, a proporção de respondentes de Bacharelado que escolheram a opção não estudou ainda a maioria desses conteúdos foi pequena.
Observa -se que, nas regiões Norte ( 3,8 \% ) e Sul ( 2,7 \% ), as proporções foram maiores do que a da média nacional ( 2,2 \% ).
Nas regiões Nordeste e Centro-Oeste, a situação foi inversa, com proporções mais baixas, correspondendo a, respectivamente, 1,5 \% e 1,6 \%.
Já na região Sudeste, a proporção foi a mesma da média nacional : 2,2 \%.
Não são observadas diferenças estatisticamente significativas entre as regiões.
Em todas as regiões, a maioria dos presentes afirmou ter estudado e aprendido muitos ou todos os conteúdos, com proporções variando entre 55,9 \%, na região Norte, a 74,1 \%, na região Sudeste.
Considerando -se separadamente as opiniões de estudantes de Licenciatura dos quatro quartos de desempenho, observa -se que, no quarto inferior, 8,5 \% ofereceram como resposta que não estudou ainda a maioria desses conteúdos, sendo 1,1 \% os do quarto superior com a mesma resposta.
As diferenças entre os alunos que optaram por esse motivo de dificuldade entre os quartos de desempenho são estatisticamente significativas.
Essas proporções são decrescentes com o desempenho.
Tendo -se em conta o quarto superior, 82,5 \% dos alunos de Licenciatura afirmaram ter estudado e aprendido muitos ou todos os conteúdos.
No outro extremo, no primeiro quarto, 49,0 \% optaram pelas mesmas categorias.
Essas proporções são crescentes com o desempenho.
Considerando -se, separadamente, as opiniões de estudantes de Bacharelado dos quatro quartos de desempenho, observa -se que, no quarto inferior, 4,5 \% ofereceram como resposta que não estudou ainda a maioria desses conteúdos, sendo 0,5 \% os do quarto superior com a mesma resposta.
A diferença entre os alunos que optaram por esse motivo de dificuldade entre os quartos de desempenho é estatisticamente significativa, com exceção da diferença entre o segundo e o terceiro quartos.
Essas proporções são decrescentes com o desempenho.
Tendo -se em conta o quarto superior, 83,3 \% dos alunos de Bacharelado afirmaram ter estudado e aprendido muitos ou todos os conteúdos.
No outro extremo, no primeiro quarto, 58,8 \% optaram pelas mesmas categorias.
Essas proporções são crescentes com o desempenho.
Na análise por Organização Acadêmica, a proporção de respondentes de Licenciatura que escolheram a opção não estudou ainda a maioria desses conteúdos foi pequena.
Observa -se que a proporção nas Universidades, com 3,5 \%, foi menor do que a média nacional ( 4,3 \% ).
Nos Centros Universitários, nas Faculdades e nos CEFET/IF, a situação foi inversa : a proporção foi mais alta ( 7,8 \%, 5,2 \% e 6,3 \%, respectivamente ).
No Gráfico 4.31a, é possível observar que as diferenças entre as Universidades em relação aos Centros Universitários e aos CEFET/IF são estatisticamente significativas.
Em todos os tipos de Organização Acadêmica, a maioria dos presentes afirmou ter estudado e aprendido muitos ou todos os conteúdos, com proporções variando entre 59,2 \%, nos CEFET/IF, a 68,3 \%, nas Universidades ( ver também Tabela II.16a, no Anexo IIA ).
Na análise por Organização Acadêmica, a proporção de respondentes de Bacharelado que escolheram a opção não estudou ainda a maioria desses conteúdos foi pequena.
Observa -se que, nos Centros Universitários ( 3,8 \% ) e nos CEFET/IF ( 4,2 \% ), as proporções foram maiores do que a da média nacional ( 2,2 \% ).
Nas Universidades ( 1,9 \% ) a situação foi inversa : a proporção foi mais baixa.
Já nas Faculdades ( 2,2 \% ) a proporção foi a mesma que a da média nacional.
No Gráfico 4.31b, é possível observar que somente a diferença entre Universidades e Centros Universitários é estatisticamente significativa ( ver também Tabela II.16b, no Anexo IIB ).
Em todos os tipos de Organização Acadêmica, a maioria dos presentes afirmou ter estudado e aprendido muitos ou todos os conteúdos, com proporções variando entre 66,8 \%, nos Centros Universitários, a 83,3 \%, nos CEFET/IF.
Em relação às Categorias Administrativas, observa -se que, nas IES Públicas, 3,5 \% ofereceram como resposta que não estudou ainda a maioria desses conteúdos, sendo 6,6 \% nas IES Privadas com a mesma resposta.
A diferença entre os alunos de Licenciatura que optaram por esse motivo de dificuldade nas duas categorias é estatisticamente significativa.
Considerando -se as IES Públicas, 67,6 \% dos alunos de Licenciatura afirmaram ter estudado e aprendido muitos ou todos os conteúdos, sendo de 61,5 \% a proporção dos que escolheram essa resposta nas IES Privadas ( ver Gráfico 4.32a e Tabela II.16a, no Anexo IIA ).
Em relação às Categorias Administrativas, observa -se que, nas IES Públicas, 1,6 \% ofereceu como resposta que não estudou ainda a maioria desses conteúdos, sendo 2,9 \% nas IES Privadas com a mesma resposta.
A diferença entre os alunos de Bacharelado que optaram por esse motivo de dificuldade nas duas categorias é estatisticamente significativa.
Considerando -se as IES Públicas, 74,8 \% dos alunos de Bacharelado afirmaram ter estudado e aprendido muitos ou todos os conteúdos.
Na outra categoria, 69,4 \% optaram pelas mesmas categorias ( ver Gráfico 4.32b e Tabela II.16b, no Anexo IIB ). 4.7 TEMPO GASTO PARA CONCLUIR A PROVA Ao responderem sobre o tempo de conclusão da prova ( Questão 9 ), cerca de quatro quintos dos estudantes de Licenciatura ( 79,9 \% ) afirmaram ter gasto Entre duas e quatro horas ( Gráfico 4.33a a Gráfico 4.36a e, no Anexo IIA, as Tabelas II.17a e II.18a ).
Considerando -se as cinco Grandes Regiões brasileiras, os percentuais dos que utilizaram Entre duas e quatro horas para finalizar a prova nas regiões Sudeste ( 78,7 \% ) e Sul ( 78,0 \% ) foram inferiores ao percentual nacional ( 79,9 \% ).
Nas demais Grandes Regiões, os percentuais de alunos de Licenciatura que dispensaram Entre duas e quatro horas para concluir a prova ficaram próximos ; 81,4 \% na região Norte ; 81,2 \%, na região Sudeste ; e 82,4 \%, na região Centro-Oeste, como mostrado no Gráfico 4.33a.
Há diferenças estatisticamente significativas entre as regiões Sudeste e Sul em relação às regiões Nordeste e Centro-Oeste, mas não entre si.
Ao responderem sobre o tempo de conclusão da prova ( Questão 9 ), quatro quintos dos estudantes de Bacharelado ( 80,0 \% ) afirmaram ter gasto Entre duas e quatro horas ( Gráfico 4.33b a Gráfico 4.36b e, no Anexo IIB, as Tabelas II.17b e II.18b ).
Considerando -se as cinco Grandes Regiões brasileiras, os percentuais dos que utilizaram Entre duas e quatro horas para finalizar a prova nas regiões Norte ( 76,6 \% ), Nordeste ( 79,2 \% ), Sul ( 78,7 \% ) e Centro-Oeste ( 77,3 \% ) foram inferiores ao percentual nacional.
Na região Sudeste ( 81,1 \% ), o percentual de alunos de Bacharelado que dispensaram Entre duas e quatro horas para concluir a prova foi superior à média nacional, como mostrado no Gráfico 4.33b.
A diferença entre as regiões não é estatisticamente significativa.
Uma vez analisadas as alternativas escolhidas pelos estudantes de Licenciatura que se situam nos diferentes quartos de desempenho, observa -se uma tendência crescente : uma proporção maior de participantes no quarto superior declarou ter gasto Entre duas e quatro horas para concluir a prova quando comparada com a proporção de participantes nos quartos inferiores.
Há diferenças estatisticamente significativas entre os quartos de desempenho, com exceção da diferença entre o segundo e o terceiro quarto de desempenho.
Uma vez analisadas as alternativas escolhidas pelos estudantes de Bacharelado que se situam nos diferentes quartos de desempenho, observa -se uma tendência crescente : uma proporção maior de participantes no quarto superior declarou ter gasto Entre duas e quatro horas para concluir a prova quando comparada com a proporção de participantes nos quartos inferiores.
Há diferenças estatisticamente significativas entre os quartos de desempenho, com exceção da diferença entre os quartos superiores de desempenho.
Na análise por Organização Acadêmica, os percentuais dos que utilizaram Entre duas e quatro horas para finalizar a prova nas Universidades ( 79,7 \% ) e Centros Universitários ( 76,6 \% ) foram inferiores ao percentual nacional ( 79,9 \% ).
Nas Faculdades ( 85,7 \% ) e nos CEFET/IF ( 81,7 \% ), o percentual foi superior ao percentual nacional, e, como mostrado no Gráfico 4.35a, há diferença estatisticamente significativa entre as Universidades e as Faculdades, e entre os Centros Universitário em relação às Faculdades e aos CEFET/IF.
Na análise por Organização Acadêmica, os percentuais dos que utilizaram Entre duas e quatro horas para finalizar a prova nas Universidades ( 79,9 \% ) e nos Centros Universitários ( 79,8 \% ) foram inferiores aos do percentual nacional.
Nos demais tipos de Organização Acadêmica, o percentual de alunos de Bacharelado que dispensaram Entre duas e quatro horas para concluir a prova ficou em 82,5 \%, nas Faculdades, e em 81,4 \%, nos CEFET/IF, conforme mostrados no Gráfico 4.35b.
As diferenças entre as Organizações Acadêmicas não são estatisticamente significativas.
Em relação às Categorias Administrativas, observa -se que, nas IES Públicas, 80,7 \% declararam ter gasto Entre duas e quatro horas para concluir a prova, sendo 77,6 \% nas IES Privadas com a mesma resposta.
A diferença entre os alunos de Licenciatura que optaram por esse motivo de dificuldade nas duas categorias é estatisticamente significativa ( ver Gráfico 4.36a e Tabela II.18a, no Anexo IIA ).
Em relação às Categorias Administrativas, observa -se que, nas IES Públicas, 80,0 \% declararam ter gasto Entre duas e quatro horas para concluir a prova, sendo, também, 79,9 \% nas IES Privadas com a mesma resposta.
A diferença entre os alunos de Bacharelado que optaram por esse motivo de dificuldade nas duas categorias não é estatisticamente significativa ( ver Gráfico 4.36b e Tabela II.18b, no Anexo IIB ).

% matched lemmas: 
\end{textsample}
